
\documentclass{article}
\usepackage{colortbl}
\usepackage{makecell}
\usepackage{multirow}
\usepackage{supertabular}

\begin{document}

\newcounter{utterance}

\centering \large Interaction Transcript for game `cladder', experiment `full\_v1.5\_default', episode 8897 with score{-}Qwen3.5{-}9B{-}sp.
\vspace{24pt}

{ \footnotesize  \setcounter{utterance}{1}
\setlength{\tabcolsep}{0pt}
\begin{supertabular}{c@{$\;$}|p{.15\linewidth}@{}p{.15\linewidth}p{.15\linewidth}p{.15\linewidth}p{.15\linewidth}p{.15\linewidth}}
    \# & \multicolumn{2}{c}{Player} && \multicolumn{2}{c}{Game Master} \\
    \hline

    \theutterance \stepcounter{utterance}  
    & & & \multicolumn{4}{p{0.6\linewidth}}{
        \cellcolor[rgb]{0.9,0.9,0.9}{
            \makecell[{{p{\linewidth}}}]{
                \texttt{\tiny{[P1$\langle$GM]}}
                \texttt{You are an expert in causal inference. The following question is not a typical commonsense query, but rather a meticulously designed question created by a professor specializing in causal inference, intended to assess the students' mastery of the course content.} \\
\texttt{Imagine a self{-}contained, hypothetical world with only the following conditions, and without any unmentioned factors or causal relationships: Yupt has a direct effect on muvq. Zupj has a direct effect on muvq. Method 1: We look directly at how yupt correlates with muvq in general. Method 2: We look at this correlation case by case according to zupj. To understand how yupt affects muvq, is it more correct to use the Method 1 than Method 2?} \\
\texttt{Start your answer with 'Yes' or 'No', followed by additional reasoning or evidence to support your explanation.} \\
            }
        }
    }
     \\ \\

    \theutterance \stepcounter{utterance}  
    & \multicolumn{4}{p{0.6\linewidth}}{
        \cellcolor[rgb]{0.561,0.737,0.561}{
            \makecell[{{p{\linewidth}}}]{
                \texttt{\tiny{[P1$\rangle$GM]}}
                \texttt{No.} \\
\\ 
\texttt{In this hypothetical world, we have a causal structure where both \$Y\_\{upt\}\$ and \$Z\_\{upj\}\$ are direct causes of \$M\_\{uvq\}\$. This implies that \$Z\_\{upj\}\$ acts as a confounder (or more precisely, a common cause in the context of the exposure{-}outcome relationship if we view \$Z\$ as a third variable influencing the outcome alongside the exposure, though here they are parallel causes).} \\
\\ 
\texttt{Method 1 calculates the marginal correlation between \$Y\_\{upt\}\$ and \$M\_\{uvq\}\$. Because \$Z\_\{upj\}\$ is a direct cause of \$M\_\{uvq\}\$, the value of \$M\_\{uvq\}\$ varies based on \$Z\_\{upj\}\$. Unless \$Z\_\{upj\}\$ is independent of \$Y\_\{upt\}\$, the variation in \$M\_\{uvq\}\$ induced by \$Z\_\{upj\}\$ will be mixed with the variation induced by \$Y\_\{upt\}\$. This results in **confounding bias**. The observed correlation in Method 1 will generally not equal the true causal effect of \$Y\_\{upt\}\$ on \$M\_\{uvq\}\$; it will be a mixture of the direct effect of \$Y\_\{upt\}\$ and the spurious association introduced by the uncontrolled variation in \$Z\_\{upj\}\$.} \\
\\ 
\texttt{Method 2 examines the correlation "case by case according to \$Z\_\{upj\}\$." This describes **stratification** (or} \\
            }
        }
    }
    & & \\ \\

    \theutterance \stepcounter{utterance}  
    & & & \multicolumn{2}{p{0.3\linewidth}}{
        \cellcolor[rgb]{0.9,0.9,0.9}{
            \makecell[{{p{\linewidth}}}]{
                \texttt{\tiny{[GM$|$GM]}}
                \texttt{No} \\
            }
        }
    }
    & & \\ \\

    \theutterance \stepcounter{utterance}  
    & & & \multicolumn{2}{p{0.3\linewidth}}{
        \cellcolor[rgb]{0.9,0.9,0.9}{
            \makecell[{{p{\linewidth}}}]{
                \texttt{\tiny{[GM$|$GM]}}
                \texttt{yes} \\
            }
        }
    }
    & & \\ \\

    \theutterance \stepcounter{utterance}  
    & & & \multicolumn{2}{p{0.3\linewidth}}{
        \cellcolor[rgb]{0.9,0.9,0.9}{
            \makecell[{{p{\linewidth}}}]{
                \texttt{\tiny{[GM$|$GM]}}
                \texttt{game\_result = LOSE} \\
            }
        }
    }
    & & \\ \\

\end{supertabular}
}

\end{document}
